\begin{figure}[ht]
\centering
\begin{tikzpicture}[x=1.0cm, y=1.0cm, font=\small]

% Left: the grid with five-point stencil.
\begin{scope}[shift={(0,0)}]
% 5x5 grid of nodes
\foreach \i in {0,...,4} {
  \foreach \j in {0,...,4} {
    \fill[enggray!55] (\i*0.75, \j*0.75) circle (1.3pt);
  }
}
% highlight centre node and its four neighbours
\def\cx{1.5} \def\cy{1.5}
\draw[thick, engblue] (\cx,\cy) -- ({\cx+0.75},\cy);
\draw[thick, engblue] (\cx,\cy) -- ({\cx-0.75},\cy);
\draw[thick, engblue] (\cx,\cy) -- (\cx,{\cy+0.75});
\draw[thick, engblue] (\cx,\cy) -- (\cx,{\cy-0.75});
\fill[engred] (\cx,\cy) circle (2.6pt);
\node[engred, font=\scriptsize, anchor=north east] at (\cx,\cy) {$4$};
\foreach \dx/\dy in {0.75/0, -0.75/0, 0/0.75, 0/-0.75} {
  \fill[engblue] ({\cx+\dx},{\cy+\dy}) circle (2.3pt);
  \node[engblue, font=\scriptsize] at ({\cx+\dx*1.35},{\cy+\dy*1.35}) {$-1$};
}
\node[anchor=north, font=\scriptsize] at (1.5, -0.4) {five-point stencil};
\end{scope}

% Right: Kronecker structure of the operator.
\begin{scope}[shift={(5.6,0.2)}]
\node[font=\small] at (2.6, 2.6)
  {$\mat{L}_{2\mathrm{D}} = \mat{T} \otimes \mat{I} + \mat{I} \otimes \mat{T}$};
\node[font=\scriptsize, text width=5.4cm, anchor=north] at (2.6, 2.2) {
  with $\mat{T}$ the $1$D second-difference matrix
  $\mathrm{tridiag}(-1, 2, -1)$. The block matrix is banded and
  symmetric positive definite; its eigenvalues are the pairwise sums
  $\lambda_{jk} = \mu_{j} + \mu_{k}$ of the $1$D spectrum, so its
  condition number is read off the $1$D problem at no extra cost.
};
% sketch of block-banded pattern
\begin{scope}[shift={(0.4,-1.9)}]
\draw[thick] (0,0) rectangle (4.2,1.4);
\foreach \k in {0,1,2} {
  \fill[engblue!35] ({\k*1.4},{1.4-0.47*(\k+1)}) rectangle ({\k*1.4+1.4},{1.4-0.47*\k});
}
% off-diagonal identity blocks (super- and sub-diagonal)
\foreach \k in {0,1} {
  \fill[engamber!30] ({\k*1.4+1.4},{1.4-0.47*(\k+1)}) rectangle ({\k*1.4+1.4+1.4},{1.4-0.47*\k});
  \fill[engamber!30] ({\k*1.4},{1.4-0.47*(\k+2)}) rectangle ({\k*1.4+1.4},{1.4-0.47*(\k+1)});
}
\node[font=\scriptsize, anchor=north] at (2.1,-0.15) {block-tridiagonal};
\end{scope}
\end{scope}

\node[anchor=north, font=\scriptsize, text width=11.5cm, align=center]
  at (4.5, -2.5) {
  The discrete Laplacian on a square grid couples each interior node to
  its four neighbours (left). Written with the Kronecker product
  (right), the two-dimensional operator is a Kronecker sum of the
  one-dimensional second-difference matrix, which makes its bandedness,
  its symmetry, and its whole spectrum inheritable from the cheap
  one-dimensional problem.
};

\end{tikzpicture}
\caption[The discrete Poisson operator and its Kronecker structure]{The
five-point finite-difference Laplacian on a structured grid and its
Kronecker-sum form. Each interior grid node contributes one row that
weights the node by $+4$ and its four neighbours by $-1$ (left),
producing a large sparse matrix that is block-tridiagonal in the natural
node ordering. The Kronecker identity $\mat{L}_{2\mathrm{D}} = \mat{T}
\otimes \mat{I} + \mat{I} \otimes \mat{T}$ (right) exposes the structure
the solver exploits: the operator is symmetric positive definite, its
bandwidth is the grid width, and its eigenvalues are sums of the
one-dimensional eigenvalues, so the conditioning of an $N^{2} \times
N^{2}$ grid problem follows from an $N \times N$ analysis. This is the
algebraic reason fast Poisson solvers and multigrid methods exist.}
\label{fig:vol02:ch09:poisson-stencil}
\end{figure}
